% ============================================================================
%  Resolución consolidada de observaciones (Artículo 46° y 47.1°, FEP01.26)
%  para la Oferta Final (\Instancia{final}).
%  Consolida el 100 % de las observaciones de la Revisión Final del Informe 2
%  (60 observaciones de Revisión 1 + 6 observaciones complementarias de
%  Revisión 2 = 66 observaciones). El Informe 3 económico no recibió informe
%  escrito de observaciones; su retroalimentación oral fue incorporada en la
%  versión económica 3.2 (Bitácora A-6 v3.2).
% ============================================================================
\begin{tablaObservaciones}{Informe 2 (Revisión Final) e Informe 3}

\grupoObservaciones{Presentación y control transversal}
\observacion[sub=0]{GEN-01}{Las tablas de respuesta ocupan una sección que altera el índice obligatorio de los subdocumentos. (rev. final, p. 3)}{La tabla de respuestas se consolida en este anexo independiente del Art. 46°. Los subdocumentos S1, S2 y S9 conservan intacta la estructura de títulos literales de la Aclaración N.° 2.}{Anexo Art. 46° y aperturas de S1, S2 y S9. Estado: Incorporada}
\observacion[sub=0]{GEN-02}{Los folios reinician por archivo. La revisión exige numeración correlativa del conjunto para la oferta final. (rev. final, pp. 1--3, 14, 48)}{Configuración de instancia actualizada a \texttt{final} con \texttt{ModoFolio\{continuo\}} para el ensamblaje definitivo del Sobre N.° 2 y conciliación con la hoja resumen de cada sobre.}{configuracion/instancia.tex y hojas resumen de sobres finales. Estado: En cierre de ensamblaje final}
\observacion[sub=0]{GEN-03}{S1, S2 y T-6 reiteran cláusulas defensivas. S1 mantiene prosa institucional de folleto. (rev. final, pp. 3--5, 7--8, 48)}{Se depuró la presentación de capacidades, compromisos y supuestos. S1 retiró misión, visión y valores y eliminó promesas genéricas no respaldadas en v3.2. T-6 remite al análisis de capacidad de S1 y al Sobre N.° 1.}{S1 §§1.1 y 1.4, S2 §§2.1--2.5 y T-6. Estado: Incorporada}
\observacion[sub=0]{GEN-04}{Las declaraciones de S3, S4 y S6 tienen seis columnas y la revisión las observa por el límite general del cuerpo. (rev. final, pp. 11, 14, 23, 48)}{Se conservan los seis campos expresamente exigidos por el Comunicado 10 §7.2 distinguiendo la declaración obligatoria de las tablas analíticas del cuerpo, con criterio uniforme en S3, S4 y S6.}{Declaraciones de uso de IA de S3, S4 y S6. Estado: Criterio documentado (pendiente validación en S3/S4/S6)}
\observacion[sub=0]{GEN-05}{S6, S8 y S9 declaran uso alto de IA con revisión concentrada. Se solicita especial atención a la simulación de riesgos y al diseño de pruebas. (rev. final, pp. 23, 31--32)}{Se documentó la revisión y trazabilidad de S9/T-13/T-17 y se actualizaron las declaraciones de S1, S2 y S9. Queda por verificar el cierre equivalente en S6 y S8.}{Declaraciones de IA de S1, S2 y S9 (incorporado); S6 y S8 (pendiente de bloque). Estado: Parcial}
\observacion[sub=0]{GEN-06}{El T-11 carece de declaración propia y los T-14, T-15 y T-18 remiten a la del capítulo. (rev. final, pp. 17--18, 27, 48)}{Individualizar la declaración de uso de IA en T-11, T-14, T-15 y T-18 al cerrar los bloques de S4 y S7, manteniendo la consolidación en A-6.}{T-11, T-14, T-15, T-18 y declaraciones de S4 y S7. Estado: Pendiente (bloque S4/S7)}

\grupoObservaciones{Empresa y comprensión de la necesidad (S1, T-6 y S2)}
\observacion[sub=1]{EMP-01}{La experiencia de mayor escala declarada cubre 120 tractocamiones frente a 374 camiones y 96.000 viajes anuales del caso, sin explicar cómo se compensa la brecha. (rev. final, pp. 4--5, 39)}{S1 §1.2 y §1.4 acreditan la compensación de escala mediante el dimensionamiento real alineado con v3.2: 37,0 FTE base en M13--M17, refuerzos de puesta en producción hasta 38,50 FTE en H6/H7, pico contratado de 39,23 FTE en M7, Mesa de Ayuda N1 6+1 FTE y servicio de campo externo por visita/rondas en los cinco terminales.}{S1 §§1.2 y 1.4 y Formulario T-6. Estado: Incorporada}
\observacion[sub=1]{EMP-02}{El gerente general es también jefe de proyecto y representante legal con dedicación de 100 \% entre M1 y M20. Falta explicar la conducción de la empresa en ese período. (rev. final, p. 5)}{Se documentaron las subrogancias gerenciales y la cobertura de funciones corporativas durante la dedicación contractual de M1--M20.}{S1 §1.3 y tabla de subrogancias. Estado: Incorporada}
\observacion[sub=1]{EMP-03}{Se vinculan entidades certificadoras y auditoras reales con expedientes ficticios. (rev. final, pp. 5, 48)}{Se armonizaron las referencias e identificadores en S1 §§1.3--1.5 y Anexo 1.A (actualizado al 25 de noviembre de 2026), remitiendo el cotejo documental al Sobre N.° 1 conforme al Art. 34.1 y Art. 39.2.}{S1 §§1.3--1.5 y Anexo 1.A (f. 7). Estado: Incorporada}
\observacion[sub=1]{EMP-04}{Las alianzas exclusivas se remiten a un capítulo y un formulario ausentes del Informe 2. (rev. final, p. 5)}{S1 §1.6 delimita explícitamente las alianzas corporativas vigentes alineadas 1:1 con v3.2 (Microsoft Azure Chile Central/Brazil South CSP, Dell, Fortinet, iWave, Teltonika, Technoton, Ground Control/Iridium, GitLab y SonarSource), radicando las cartas de respaldo en el Sobre N.° 1 y su resumen en el Capítulo 12 / Formulario T-8.}{S1 §1.6 y Formulario T-8. Estado: Incorporada}
\observacion[sub=2]{NEC-01}{La estimación de 338,15 viajes diarios de temporada alta y la cota de 450 viajes diarios dependen de un supuesto de 210 viajes diarios en valle. (rev. final, pp. 7--8)}{S2 documenta la ecuación de balance anual, la cota de estrés y su calibración obligatoria con el histórico real del sistema en H1 (M2).}{S2 §§2.1 y 2.2.3. Estado: Incorporada}
\observacion[sub=2]{NEC-02}{El diagnóstico presenta los umbrales de continuidad como compromisos propios y usa la expresión «análisis pericial». (rev. final, pp. 7--8)}{S2 atribuye expresamente a las Bases Oficiales la disponibilidad (99,9 \%), RTO ($\le 4$ h), RPO ($\le 15$ min) y autonomía (72 h), corrigiendo las citas de numerales y retirando la expresión «análisis pericial».}{S2 §§2.2.2 y 2.5.1. Estado: Incorporada}
\observacion[sub=2]{NEC-03}{La figura de cadena de valor queda aislada por espacios en blanco. (rev. final, p. 7)}{La Figura 2.1 fue redibujada en TikZ vectorial nativo a 9,5 pt integrada al flujo de S2 §2.2, recompilada y verificada sin espacios muertos ni errores en el PDF final (30 págs.).}{S2 §2.2, Figura 2.1 (PDF final verificado). Estado: Incorporada}

\grupoObservaciones{Alcance y cumplimiento técnico (S3 y T-12)}
\observacion[sub=3]{ALC-01}{SA-01 fija febrero de 2027 como inicio y el calendario depende de esa fecha para la ventana de M16. (rev. final, pp. 11--12, 46--47)}{Separar los meses contractuales del calendario civil y evaluar inicios alternativos contra las ventanas de montaje y puesta en producción, sin supeditar la oferta a SA-01 ni desplazar hitos E-25.}{S3 supuestos y calendario, S7 §§7.2--7.3, T-14/T-15/T-18. Estado: Pendiente (bloque S3/S7)}
\observacion[sub=3]{ALC-02}{La reversión se resume en una afirmación de que no se pierde información. (rev. final, p. 11)}{CP-SYS-20 en T-17 y el procedimiento de 6 pasos de T-18 articulan disparadores, buffer de 72 h y conciliación; resta verificar el enlace explícito en el cierre de S3.}{S3 §3.3, T-17 (CP-SYS-20) y T-18. Estado: Parcial (T-17 cerrado; enlace en S3 pendiente)}
\observacion[sub=3]{ALC-03}{El T-12 contiene 42 requerimientos propios y usa los RT como referencias, sin declarar su cumplimiento individual. (rev. final, pp. 11--12, 33, 46--47)}{Inventariar cada requisito del caso y cada RT transversal en T-12 con cumplimiento individual, componente, sección y folio de evidencia, conciliado con T-17.}{T-12, S3 trazabilidad y T-17. Estado: Pendiente (bloque S3/T-12)}
\observacion[sub=3]{ALC-04}{El T-12 declara «Sí» uniforme, incluso en RF-003, RNF-005 y RF-014, que dependen de autorizaciones o pruebas futuras. (rev. final, pp. 11--12, 46--47)}{Distinguir compromiso contractual, evidencia disponible y verificación futura en las fichas RF-003, RNF-005 y RF-014, explicitando dependencias y plan de resolución.}{T-12 fichas RF-003, RNF-005 y RF-014, S3 y T-16. Estado: Pendiente (bloque S3/T-12)}
\observacion[sub=3]{ALC-05}{La numeración de las tablas 3.1--3.7 del T-12 coincide con la del S3. (rev. final, pp. 11--12, 47)}{Asignar prefijo propio al Formulario T-12 y actualizar todas las remisiones cruzadas en S3 y T-12.}{T-12 y S3. Estado: Pendiente (bloque S3/T-12)}
\observacion[sub=3]{ALC-06}{El cuerpo del S3 se refiere a versiones previas del informe. (rev. final, pp. 11--12)}{Redactar S3 como propuesta autónoma y concentrar el historial de correcciones exclusivamente en este anexo del Art. 46°.}{S3 y anexo Art. 46°. Estado: Pendiente (bloque S3)}

\grupoObservaciones{Arquitectura y datos (S4, T-11 y S5)}
\observacion[sub=4]{ARQ-01}{Flutter, React y Next.js no tienen versiones y soporte definidos. La versión Go declarada carece de respaldo verificable para la fecha de oferta. (rev. final, pp. 13--14, 43)}{Fijar versiones verificadas a noviembre de 2026 con fuente de soporte, fin de vida y política de actualización, alineando S4 con T-11.}{S4 §4.1 y T-11. Estado: Pendiente (bloque S4)}
\observacion[sub=4]{ARQ-02}{Los ocho ADR llevan una misma fecha y carecen de estado por decisión. (rev. final, pp. 14--15)}{Registrar fecha efectiva, estado, alternativa seleccionada, descartes y criterio para cada uno de los ocho ADR.}{S4 §4.1, Tabla 4.2 y registro ADR. Estado: Pendiente (bloque S4)}
\observacion[sub=4]{ARQ-03}{La mesa de ayuda y la dotación operativa se remiten al S10 sin dimensionamiento. (rev. final, pp. 15--16, 18, 44)}{Incorporar en S4/S10 el dimensionamiento ya cuadrado en v3.2 y S1 (Mesa N1 24×7 de 6 agentes regulares + 1 FTE suplencias = 7 FTE; campo 0,5 FTE interno + visitas/rondas externas).}{S4 §4.2, S10--S11 y T-15 (cuadrado con S1 y v3.2). Estado: Pendiente (bloque S4/S10)}
\observacion[sub=4]{ARQ-04}{El alcance de ExpressRoute hasta Chile Central queda sujeto a confirmación del proveedor. (rev. final, pp. 17--18, 29, 31, 46--47)}{Documentar factibilidad del circuito ExpressRoute con operadores evaluados, definir conectividad alternativa conforme a Bases y vincular su exposición en T-16.}{S4 §4.2, T-11, S8/T-16 y T-17. Estado: Pendiente (bloque S4/S8)}
\observacion[sub=4]{ARQ-05}{El DR de datos personales en Brazil South carece de vía habilitante documentada. R-24 depende de una actuación del mandante antes de H2 y afecta H7. (rev. final, pp. 17--20, 30--31, 45--47)}{Documentar la vía habilitante aplicable bajo Ley N.° 21.719 y definir mecanismo alternativo admisible con evidencia de RTO/RPO si no se obtiene autorización, sin presuponer aprobación del mandante.}{S4 §4.2 DR, S5 §5.1.6, S8/T-16 R-24 y T-17. Estado: Pendiente (bloque S4/S5/S8)}
\observacion[sub=4]{ARQ-06}{La tabla del T-11 se identifica con el mismo número que la tabla de tecnologías de S4 y remite las referencias al capítulo. (rev. final, pp. 17--18, 47)}{Asignar numeración propia al Formulario T-11 y localizar sus fuentes técnicas directamente en el formulario.}{T-11 y S4. Estado: Pendiente (bloque S4/T-11)}
\observacion[sub=4]{ARQ-07}{Persisten brechas de cobertura real, parque de terceros, conductor sin rastro, consultas sin respuesta y autonomía de la solución satelital referida en la revisión. (rev. final, p. 18)}{Explicitación de dato faltante, prueba o consulta, plazo de resolución y contingencia por cada brecha, separando la autonomía ofertada de antecedentes previos.}{S4 §§4.1--4.2, S3, T-16 y T-17. Estado: Pendiente (bloque S4/S3)}
\observacion[sub=4]{ARQ-08}{[Observación complementaria Rev. 2] La Figura 4.13 dibuja la «Autoridad tributaria» dentro del bloque de San Bernardo, junto al sistema contable, sin mostrar el enlace externo con el emisor ni con la capa anticorrupción. (rev. final, p. 18)}{Reubicar la entidad «Autoridad tributaria» fuera del perímetro físico de San Bernardo en la Figura 4.13 y graficar explícitamente el enlace externo seguro con el emisor tributario y la capa anticorrupción.}{S4 §4.2, Figura 4.13. Estado: Pendiente (bloque S4)}
\observacion[sub=5]{DAT-01}{El diccionario de datos se incorporó dentro del archivo del S5. (rev. final, pp. 20, 48)}{Publicar el diccionario de datos como anexo independiente (\texttt{AUDIT-Subdocumento5-Anexos.pdf}) y mantener en el cuerpo de S5 los modelos, reglas y remisión al anexo.}{S5 y AUDIT-Subdocumento5-Anexos. Estado: Pendiente (bloque S5)}
\observacion[sub=5]{DAT-02}{Las tablas 5.2, 5.3, 5.5 y 5.6 contienen prosa explicativa en las celdas. (rev. final, pp. 20, 48)}{Trasladar reglas extensas y excepciones al cuerpo o diccionario, reservando las celdas para códigos y atributos comparables.}{S5, Tablas 5.2, 5.3, 5.5 y 5.6. Estado: Pendiente (bloque S5)}
\observacion[sub=5]{DAT-03}{La declaración de S5 deja una página de continuación con encabezado y sin filas. (rev. final, pp. 20, 48)}{Eliminar el salto vacío y verificar el cierre, encabezados y folios al recompilar S5.}{S5, Declaración de uso de IA. Estado: Pendiente (bloque S5)}
\observacion[sub=5]{DAT-04}{El 99,9 \% mensual de servicios críticos en S5 debe cotejarse con el SLA de S4. (rev. final, p. 20)}{Conciliar meta, servicios cubiertos, ventana y fórmula de disponibilidad entre S5, S4, S10 y T-17.}{S5 Tabla 5.4, S4, S10 y T-17. Estado: Pendiente (bloque S5/S4)}
\observacion[sub=5]{DAT-05}{[Observación complementaria Rev. 2] La entidad TRANSPORTISTA cumple tres papeles en la Figura 5.1 con el mismo \texttt{id\_transportista} sin regla cuando difieren, y LIQUIDACION\_DETALLE tiene PK por viaje y no por asignación, con riesgo de doble liquidación ante reasignación. (rev. final, p. 20)}{Desacoplar roles de titularidad/empleador/responsable de liquidación en el modelo conceptual y referenciar \texttt{id\_asignacion} en \texttt{LIQUIDACION\_DETALLE} con restricción de unicidad que impida liquidar dos veces el mismo tramo de un viaje.}{S5 Figura 5.1 y Tabla 5.33. Estado: Pendiente (bloque S5)}

\grupoObservaciones{Metodologías (S6, T-9 y T-10)}
\observacion[sub=6]{MET-01}{El plan de comunicaciones cubre cinco grupos y no articula los trece actores de S2 ni define destinatario, frecuencia o canal. (rev. final, pp. 21, 23)}{Construir matriz de comunicaciones para los trece actores identificados en S2 con responsable, destinatario, frecuencia, canal y escalamiento.}{S6 §6.1, T-9 y S2 §2.4. Estado: Pendiente (bloque S6)}
\observacion[sub=6]{MET-02}{Falta identificar quién prioriza y acepta incrementos por Curimón y quién lo reemplaza. (rev. final, pp. 21, 23)}{Definir rol titular y suplente del mandante, atribuciones y plazos de decisión, diferenciando aceptación de incremento de aceptación contractual (H4/H9).}{S6 §§6.1--6.2, T-9/T-10 y S7. Estado: Pendiente (bloque S6)}
\observacion[sub=6]{MET-03}{T-9 y T-10 se limitan a remisiones, con poca información utilizable por sí mismos. (rev. final, pp. 21--22)}{Incorporar contenido autosuficiente de gestión en T-9 y de desarrollo en T-10 conforme a sus instrucciones oficiales.}{T-9 y T-10. Estado: Pendiente (bloque S6)}
\observacion[sub=6]{MET-04}{Los ocho sprints de A09 y cinco de A20 no tienen un plan de entregas por contexto o requerimiento. Sólo se detalla G13 en el sprint 7. (rev. final, pp. 22--23, 26--27, 48)}{Asignar por sprint los requisitos, incrementos, dependencias, capacidad neta, prueba, usuario revisor y criterio de salida, conciliado con T-15 y T-17.}{S6 §6.2, S7, T-10, T-12, T-15 y T-17. Estado: Pendiente (bloque S6/S7)}
\observacion[sub=6]{MET-05}{La deuda técnica se descuenta en sprints, pero carece de paquete y horas durante la operación Kanban. (rev. final, pp. 22--23)}{Incorporar presupuesto de capacidad y paquete de deuda técnica para M21--M56 en el modelo de operación Kanban alineado con v3.2.}{S6 §6.2, S7 §7.2, T-15 y S11. Estado: Pendiente (bloque S6/S11)}
\observacion[sub=6]{MET-06}{La viñeta de ramas temáticas de GitFlow cita normas ISO que no sustentan esa práctica. (rev. final, p. 23)}{Retirar las citas ISO improcedentes en control de versiones y respaldar el flujo de ramas con referencia técnica pertinente.}{S6 §6.2 y referencias. Estado: Pendiente (bloque S6)}
\observacion[sub=6]{MET-07}{Las tablas 6.1, 6.2, 6.3 y 6.6 contienen prosa. La Tabla 6.2 utiliza una descripción arquitectónica que no coincide con S4. (rev. final, p. 23)}{Trasladar la prosa extensa al texto y armonizar la Tabla 6.2 con los contenedores y ADR vigentes de S4.}{S6 Tablas 6.1, 6.2, 6.3 y 6.6 y S4 §4.1. Estado: Pendiente (bloque S6)}
\observacion[sub=6]{MET-08}{[Observación complementaria Rev. 2] El S6 no menciona la adhesión de los 148 transportistas ni la trata como condición del método, aunque S7 le asigna paquetes EDT 9.1--9.3. (rev. final, p. 23)}{Incorporar el gobierno del frente de adhesión de los 148 transportistas en S6 §6.1 y T-9 (responsable, cadencia de seguimiento, umbral de avance y escalamiento ante desvío), enlazado con los paquetes 9.1--9.3 de S7.}{S6 §6.1, T-9 y S7 EDT 9.1--9.3. Estado: Pendiente (bloque S6)}

\grupoObservaciones{EDT, cronograma e implantación (S7, T-14, T-15 y T-18)}
\observacion[sub=7]{PLA-01}{El primer nivel de la EDT se justifica genéricamente, divide el software por etapa y carece de agrupación intermedia. (rev. final, pp. 24, 27, 48)}{Agrupar el software por capacidad o contexto verificable representando la etapa como atributo de planificación.}{S7 §7.1 y T-14. Estado: Parcial (bloque S7)}
\observacion[sub=7]{PLA-02}{Las láminas EDT repetían el diccionario en vez de mostrar la jerarquía. (rev. final, pp. 24, 27)}{El S7 revisado incorpora el árbol general y cinco vistas jerárquicas por rama, conservando el diccionario en T-14.}{S7 §7.1, Figuras 7.3--7.8 y T-14. Estado: Incorporada}
\observacion[sub=7]{PLA-03}{Los cinco paquetes de innovación carecían de meses concretos en T-14 y de actividad identificable en la red. (rev. final, pp. 24, 27, 37)}{T-14 y S7 fechan los cinco paquetes; resta conciliar esfuerzo y dependencias cruzadas en T-15 y T-19.}{T-14, S7, T-15 y T-19. Estado: Parcial (bloque S7/S13)}
\observacion[sub=7]{PLA-04}{El esfuerzo se calcula desde dotaciones mensuales y no por tarea y paquete. (rev. final, pp. 25, 27, 48)}{Estimar horas por entregable, tarea y rol y conciliar con las curvas de dotación de v3.2 (pico 39,23 FTE en M7 y refuerzos 38,50 FTE en H6/H7).}{S7 §7.2 y T-15. Estado: Pendiente (bloque S7)}
\observacion[sub=7]{PLA-05}{A09 y A20 agrupan toda la construcción de cada etapa sin secuenciar los contextos dependientes. (rev. final, pp. 25, 27, 48)}{Descomponer A09 y A20 enlazando vigencias, jornada, asignación, documento y costo según dependencias reales y recalcular holguras/ruta crítica.}{S7 §§7.2--7.3, T-14 y T-15. Estado: Pendiente (bloque S7)}
\observacion[sub=7]{PLA-06}{La migración histórica y el estudio de costo por ruta no tienen precedencias ni holguras en la red. (rev. final, p. 25)}{Incorporar los paquetes 7.2 y 2.5 a la red CPM con predecesoras, sucesoras, holguras y responsables.}{S7 §§7.2--7.3, T-14 y T-15. Estado: Pendiente (bloque S7)}
\observacion[sub=7]{PLA-07}{La Gantt por paquete no muestra dependencias ni la ruta crítica de manera conciliable con T-15. (rev. final, pp. 25, 27, 48)}{La vista CPM incorpora ruta crítica; resta enlazar dependencias a nivel de paquete con T-15.}{S7 §7.3.2 y T-15. Estado: Parcial (bloque S7)}
\observacion[sub=7]{PLA-08}{La ruta crítica encadena levantamiento, arquitectura, infraestructura, construcción y certificación con escaso paralelismo. (rev. final, p. 25)}{Evaluar inicios parciales por capacidad y ambiente y justificar técnicamente cada dependencia en serie.}{S7 §§7.2--7.3 y T-15. Estado: Pendiente (bloque S7)}
\observacion[sub=7]{PLA-09}{Las láminas Gantt de pocas filas y las 42 páginas gráficas fragmentaban la lectura. (rev. final, pp. 25, 27)}{S7 concentra la Gantt en panorama, vista CPM y diez vistas temáticas, separando la EDT.}{S7 §7.3.2, Figuras 7.12--7.23. Estado: Incorporada}
\observacion[sub=7]{PLA-10}{La primera marcha blanca coincide con la temporada de fruta sin declarar la disponibilidad de usuarios de torre y terminales como dependencia. (rev. final, p. 26)}{Declarar ventanas, cupos, suplencias y escalamiento con el mandante preservando los hitos y las seis condiciones del Art. 17.3.}{S7 §7.3, T-18, S8/T-16 y S9. Estado: Pendiente (bloque S7/S8)}
\observacion[sub=7]{PLA-11}{A18 combina integración telemática en M13--M15 con montaje en M16--M18, pese a la restricción estacional. (rev. final, pp. 26--27)}{Separar ambas actividades en la red T-15 con recursos, precedencias y capacidad de montaje por terminal alineada con v3.2.}{S7 §7.3.2 y T-15. Estado: Parcial (bloque S7)}
\observacion[sub=7]{PLA-12}{El piloto A11 coincide con meses de cierres de Los Libertadores. La holgura se justifica con un cierre individual de hasta 12 días. (rev. final, p. 26)}{Evaluar cierres sucesivos en invierno, consumo de holgura y alternativa de ruta/reprogramación que preserve la aceptación del piloto.}{S7 §7.3, T-15 (A11/A18) y T-16. Estado: Pendiente (bloque S7/S8)}
\observacion[sub=7]{PLA-13}{El plan de montaje no incorpora las 210 balizas de semirremolque, las 21 de reposición ni los lectores de portería declarados para la innovación 3. (rev. final, pp. 27, 48)}{Incorporar paquete, sitios, actividad, rendimiento y diferenciación entre montaje y stock de reposición (210 + 21 balizas; 182 + 19 kits a bordo), conciliado con v3.2, T-11 y T-19.}{S7 §§7.1--7.3, T-14, T-15, T-18, T-11 y T-19. Estado: Pendiente (bloque S7)}
\observacion[sub=7]{PLA-14}{[Observación complementaria Rev. 2] El Formulario T-14 no contiene la Carta Gantt (la remite a la sección 7.3.2 de S7) y la Etapa 2 carece de actividad de estabilización post-producción equivalente a A21 de Etapa 1. (rev. final, p. 27)}{Adjuntar la Carta Gantt directamente dentro del Formulario T-14 conforme a sus instrucciones e incorporar actividad formal de estabilización post-producción de Etapa 2 en M21--M24 (conciliada con el cierre de M21--M24 presupuestado en v3.2).}{Formulario T-14, S7 §7.3 y T-15. Estado: Pendiente (bloque S7)}

\grupoObservaciones{Riesgos (S8 y T-16)}
\observacion[sub=8]{RIE-01}{La mitad del registro del cliente agrupa riesgos internos de audIT. R-06 trata como incertidumbre un retiro anunciado y ya resuelto por diseño. (rev. final, pp. 28, 31, 48)}{Clasificar riesgos por origen y actuación requerida del mandante, reclasificar R-06 como antecedente del ADR de caché y mantener visibles los riesgos contractuales abiertos.}{S8 §§8.1--8.3 y T-16. Estado: Pendiente (bloque S8)}
\observacion[sub=8]{RIE-02}{La convivencia y el retiro del sistema de 2013 no aparecen en el registro de riesgos. (rev. final, pp. 29, 31, 48)}{Incorporar exposición de convivencia y retiro del TMS 2013 en M16, M21 y M24 con disparadores, conciliación, reversión y efecto sobre hitos.}{S8/T-16, S3, S7 y T-18. Estado: Pendiente (bloque S8)}
\observacion[sub=8]{RIE-03}{Faltan en el registro los riesgos de adopción de las innovaciones 1 y 2 y su riesgo residual. (rev. final, pp. 29, 31)}{Registrar eventos, disparadores, mitigación y nivel residual de adopción de las innovaciones 1 y 2, conciliados con T-19.}{S8/T-16 y T-19. Estado: Pendiente (bloque S8)}
\observacion[sub=8]{RIE-04}{El árbol de fallas de asignación a conductor no apto carece de cuantificación. (rev. final, pp. 29, 31)}{Definir protocolo de conteo y cuantificación previa al paralelo M13--M14 declarando incertidumbre y límites de estimación.}{S8 §8.2 Figura 8.3 y T-16. Estado: Pendiente (bloque S8)}
\observacion[sub=8]{RIE-05}{El análisis costo-beneficio de respuestas no presenta una razón de retorno. (rev. final, p. 30)}{Calcular costo incremental y pérdida esperada evitada manteniendo los valores monetarios exclusivamente en el Sobre N.° 3 (Art. 50.2) y el criterio técnico en S8.}{S8 §8.3, T-16 y Sobre N.° 3. Estado: Pendiente (bloque S8)}
\observacion[sub=8]{RIE-06}{R-17 se acepta con seguimiento, aunque la simulación reporta 21,6 \% de atraso de H1 y 16,1 \% de H2. (rev. final, pp. 30--31, 48)}{Definir umbrales de alerta temprana, escalamiento y contingencia para H1 y H2 y recalcular el residual sin alterar hitos contractuales.}{S8 §8.3 y T-16 (R-17). Estado: Pendiente (bloque S8)}
\observacion[sub=8]{RIE-07}{Las tablas 8.1--8.3 se repiten entre S8 y T-16 y las remisiones son ambiguas. (rev. final, p. 31)}{Asignar numeración propia al Formulario T-16 y depurar remisiones cruzadas con S8.}{S8 §8.3 y T-16. Estado: Pendiente (bloque S8)}
\observacion[sub=8]{RIE-08}{[Observación complementaria Rev. 2] R-22 trata suplantación digital pero no la desconexión física de antena/CAN/lector por conductores externos; la mitigación de R-14 es débil frente a 258 conductores externos al no apoyarse en el incentivo de la innovación 4 ni en el anexo de adhesión. (rev. final, p. 31)}{Ampliar R-22 (o crear subriesgo ligado a R-14) para desconexión física/sabotaje de telemetría con disparador de pérdida de señal en ruta, y fortalecer la mitigación de R-14 vinculándola explícitamente al comodato/incentivo de la innovación 4 y a las cláusulas del anexo de adhesión.}{S8 §8.3 y T-16 (R-14 y R-22). Estado: Pendiente (bloque S8)}

\grupoObservaciones{Calidad y aceptación (S9, T-13 y T-17)}
\observacion[sub=9]{CAL-01}{El calendario de ensayos y las puertas de calidad se muestran sólo en tablas de detalle. (rev. final, p. 32)}{S9 incorpora la vista temporal vectorial (Figura 9.2) alineada con S7/T-15 y los hitos E-25 (actualizada en versión final con A20 en M13--M18 e hitos H3, H6, H8 y H11) y su explicación operativa.}{S9 §9.3.2, Figura 9.2 (PDF final verificado). Estado: Incorporada}
\observacion[sub=9]{CAL-02}{CP-HW-08 alude a polvo minero y CP-HW-06 duplica el ensayo de alarma SOS de CP-PERF-06. (rev. final, pp. 32, 34)}{CP-HW-08 fue saneado a estanqueidad IP67 en transporte terrestre (tanto en \texttt{catalogo-T-17.tex} como en \texttt{T-17.md}); CP-HW-06 verifica entrega SOS por evento con cobertura real y CP-PERF-06 mide P99 bajo degradación reproducible ($N \ge 100$).}{Formulario T-17 (\texttt{CP-PERF-06}, \texttt{CP-HW-06}, \texttt{CP-HW-08}). Estado: Incorporada}
\observacion[sub=9]{CAL-03}{Los 29 criterios de aceptación del capítulo 18 del caso no se vinculan a casos de prueba ni a actas. (rev. final, pp. 33--34, 48)}{La matriz \texttt{9-rtm-29} en T-17 traza los 29 criterios del Capítulo 18 a requisitos, casos de prueba, hitos y responsables; se cuadró además la trazabilidad de \texttt{CP-SYS-05} (\texttt{RF-002}) y \texttt{CP-PERF-03} (\texttt{RNF-014}, \texttt{RT-07.04}, \texttt{RT-07.07}).}{Formulario T-17 (Tabla \texttt{9-rtm-29} y catálogo). Estado: Incorporada}
\observacion[sub=9]{CAL-04}{La UAT de Etapa 1 coincide con el comienzo de la temporada de fruta, sin disponibilidad acordada de usuarios de terminal y despacho. (rev. final, p. 34)}{S9 §9.3.2 y T-13 §5 establecen ventanas UAT acotadas por turno fuera de horario punta, cupos nominales por rol (torre de despacho, prevención, mantenimiento, finanzas y terminales), suplentes designados y escalamiento en 48 h ante indisponibilidad para preservar H5; además, los simulacros semestrales de DR se fijan en junio y noviembre fuera de temporada alta.}{S9 §9.3.2 y Formulario T-13 §5. Estado: Incorporada}
\observacion[sub=9]{CAL-05}{El S9 ubica A12 en M10--M12 y el T-13 la expresa como M12/H5. (rev. final, p. 34)}{S9 y T-13 distinguen la ventana de ejecución de A12 (M10--M12) del hito de certificación H5 en M12, en total cuadratura con T-15 y E-25.}{S9 Tabla 9.5 y T-13 apartado 5. Estado: Incorporada}
\observacion[sub=9]{CAL-06}{El T-17 incluye un protocolo general y carece de protocolos por cada hito y del producto final. (rev. final, pp. 34, 48)}{El Formulario T-17 formaliza 13 fichas de protocolos individuales (H1 a H12 y Producto Final) alineadas a los Arts. 18.1--18.3 y garantía de correcto funcionamiento en M21 (Art. 37).}{Formulario T-17, protocolos de hitos H1--H12 y cierre. Estado: Incorporada}
\observacion[sub=9]{CAL-07}{[Observación complementaria Rev. 2] Las pruebas de integración con plataformas GPS de terceros (CP-INT-07 y CP-INT-08) se ejecutan sólo contra dobles simulados; ningún caso de staging o UAT define la muestra autorizada real con las tres plataformas. (rev. final, p. 34)}{Se incorporó en S9 §9.3.2, T-13 §5 y \texttt{CP-UAT-09} de T-17 la validación obligatoria en staging/UAT sobre muestra autorizada real de al menos una unidad activa por cada una de las tres plataformas GPS de terceros (SUP-02) previo al cierre de H5/H10, complementando los dobles simulados de CI (\texttt{CP-INT-07} a \texttt{CP-INT-09}).}{S9 §9.3.2, T-13 §5 y Formulario T-17 (\texttt{CP-UAT-09}). Estado: Incorporada}

\grupoObservaciones{Innovaciones y valorización (S13 y T-19)}
\observacion[sub=13]{INN-01}{La programación por oportunidad de patio tiene TRL 2 y metas relativas a una línea base sin medir. (rev. final, pp. 35--37, 45)}{Definir línea base, método de medición en piloto, umbrales observables y fecha de decisión, manteniendo CLP 0 de beneficio en escenario base de v3.2.}{S13/T-19 innovación 2 y análisis económico v3.2. Estado: Pendiente (bloque S13)}
\observacion[sub=13]{INN-02}{La innovación 3 presenta como novedad el kilometraje real de semirremolques, que apoya el criterio de aceptación 26. (rev. final, pp. 36--37)}{Clasificar el kilometraje real como alcance base (criterio 26) y acotar la innovación 3 a verificación de acople y registro térmico adicional (ya alineado en S2 y v3.2).}{S13/T-19 innovación 3 (ya alineado en S2 y v3.2). Estado: Parcial (pendiente cierre en S13/T-19)}
\observacion[sub=13]{INN-03}{La innovación 5 presenta el catálogo de paradores del criterio 28 como parte del mérito innovador. (rev. final, pp. 35--37, 45)}{La ficha de prevención de fatiga en T-19 distingue el mérito innovador del cumplimiento obligatorio del criterio 28.}{Formulario T-19, ficha innovación 5. Estado: Incorporada}
\observacion[sub=13]{INN-04}{Las cinco fichas omiten el riesgo residual. La percepción de vigilancia en la innovación 5 se describe con probabilidad media a alta e impacto alto. (rev. final, pp. 36--37)}{Las cinco fichas de T-19 incorporan riesgo inherente, medidas de tratamiento y evaluación de riesgo residual mitigado; conciliar con S8/T-16.}{Formulario T-19, Fichas 1 a 5. Estado: Parcial (conciliación con S8)}
\observacion[sub=13]{INN-05}{Las innovaciones 2, 3 y 5 carecen de base y magnitud para sus beneficios, aunque el informe remite la valorización al Sobre 3. (rev. final, p. 36)}{Definir indicadores físicos medibles en piloto manteniendo en v3.2 CLP 0 de beneficio base para innovaciones 2 y 3, monetización pendiente para 1 y 5, y recupero bruto referencial separado de los ingresos de audIT para la innovación 4.}{S13/T-19 y Sobre N.° 3 (v3.2). Estado: Parcial (v3.2 cerrado; cierre narrativo en S13)}
\end{tablaObservaciones}
